\chapter{Motivazioni e obiettivi dello stage}
\label{chap:motivazioni-obiettivi}

Questo capitolo discute i problemi che l'organizzazione ospitante ha inteso affrontare con lo stage proposto, il ruolo degli stage nella strategia di innovazione dell'azienda, gli obiettivi e i vincoli dello stage, e le ragioni della scelta da parte dello studente.

\section{Il ruolo degli stage nella strategia dell'azienda}
\label{sec:ruolo-stage}

% TODO: collegamento con quanto discusso in chiusura del capitolo 1; come gli
% stage si inseriscono nel rapporto tra l'azienda e l'innovazione; continuità
% tra stage passati, presente e futuri.


\section{Il problema affrontato: l'accessibilità digitale come esigenza aziendale}
\label{sec:problema-affrontato}

% TODO: inquadramento sintetico, dal punto di vista aziendale, della crescente
% rilevanza normativa ed economica dell'accessibilità digitale (obblighi di
% legge, scadenze regolatorie, richieste di clienti o mercato) come motivazione
% concreta alla base della proposta di stage. Solo un inquadramento sintetico:
% l'approfondimento tecnico e normativo completo è nella sezione~\ref{sec:quadro-normativo}
% del capitolo~\ref{chap:sviluppo-estensione}.


\section{Obiettivi e vincoli dello stage}
\label{sec:obiettivi-vincoli}

% TODO: obiettivi assegnati allo stage e vincoli (temporali, tecnologici,
% organizzativi) posti dall'azienda, in relazione alla visione di innovazione
% descritta nella sezione precedente.


\section{La scelta dello stage}
\label{sec:scelta-stage}

% TODO: motivazione della scelta di questo stage rispetto ad altre proposte
% disponibili, dal punto di vista dello studente.
